\section{Introduction}
A finite stochastic register has a continuum of probability distributions. It is therefore not immediate that reversibility of a numerical experiment forces a finite-state implementation to be deterministic. The implementation need not respect the group relations: a physically trivial word may act nontrivially on its hidden register. Approximate agreement makes this distinction especially important. Our first result shows that, for one stationary implementation valid at every word, the transition randomness can nevertheless be removed without a loss in width or accuracy. The proof retains the physical group coordinate throughout the reduction.

Throughout the stationary part of the paper, $H$ is an arbitrary compact Hausdorff group and $\mathcal A$ is a finite nonempty alphabet with letters $u_a\in H$ generating a dense subgroup. We use the execution convention
\[
             u_{a_1\cdots a_n}=u_{a_n}\cdots u_{a_1}.
\]
There are finite nonempty sets of seeds $\mathcal S$ and queries $\mathcal J$, and continuous maps $F_{s,j}:H\to Y_j$. Each $Y_j$ is a nonempty compact convex subset of a finite-dimensional real normed space. A categorical output is the whole vector in a probability simplex, with total variation as its norm. Distinct incompatible measurements are separate queries, not a postulated joint distribution.

\begin{definition}[Stationary realization]\label{def:stationary55}
A realization on $k$ labels consists of initialization rows $\alpha_s\in\Delta_k$, row-stochastic matrices $T_a$ of order $k$, and decoder rows $D_j(i)\in Y_j$. All are independent of the word length and clock. Its error is
\[
 \sup_{s,j,w\in\mathcal A^*}
 \norm{\alpha_sT_{a_1}\cdots T_{a_n}D_j-F_{s,j}(u_w)}_j.
\]
The empty word is included. Write $M_\varepsilon$ for the least such $k$ at error at most $\varepsilon$, or $\infty$ if there is none. A permutation realization has permutation matrices for all commands; its initialization need not be a point mass.
\end{definition}
All advertised labels count, even if some are unused. A decoder vector is not an additional sampled-answer register. When an actual finite answer register is charged, its size is included separately as in Definition~\ref{def:machine54}.

Theorem~\ref{thm:purification55} proves that the minimum $M_\varepsilon$ is unchanged by restricting to permutation realizations. Theorem~\ref{thm:stationaryminimum55} expresses it by finitely many permutation choices at each $k$, followed by a compact feasibility problem. In particular, the transition matrices need not be searched over a continuum. Random initialization remains essential: Proposition~\ref{prop:C655} has unrestricted stationary minimum five and deterministic component-quotient minimum six.

For an open normal subgroup $U$ of $H$ put
\begin{equation}\label{eq:quotientradius55}
 r(U)=\max_{s,j,g\in H}\rad_{Y_j}\bigl(F_{s,j}(gU)\bigr),
 \qquad
 \rad_Y(A)=\min_{y\in Y}\sup_{x\in A}\norm{x-y}.
\end{equation}
These maxima and minima are attained: continuous images of compact cosets are compact, and the radius depends continuously on their translate. The finite-quotient criterion, Theorem~\ref{thm:finitequotient55}, says precisely
\[
          M_\varepsilon<\infty
          \quad\Longleftrightarrow\quad
          r(U)\le\varepsilon\text{ for some open normal }U.
\]
This is an equality-inclusive criterion, valid without an executable return. If $H/H^\circ$ is finite, it reduces to the largest constrained radius of a connected-component image.

A different resource is the least width $W_{N,\varepsilon}$ when all command rows may depend on a free clock and the entire machine may be redesigned at horizon $N$. The stationary purification theorem does not identify this resource with $M_\varepsilon$. The distinction is necessary: a single irrational-rotation letter admits a one-label clocked realization at each horizon. For clocked realizations we prove a converse by localized conditional-moment contraction. With cofinite exact returns and finitely many components, it gives the sharp component-radius theorem, including boundary attainment and an occupation bound for all cuts of any fixed width. Section~\ref{sec:phases55} replaces cofiniteness by the existence of a nonempty exact return, with the necessary refinement by length phase. Theorem~\ref{thm:profinite55} treats arbitrary component groups on both strict sides of the radius, while Theorem~\ref{thm:zero55} identifies the zero-error boundary with finite observable quotients. A concrete $2$-adic interface has $W_{N,0}=\Theta(N)$ and an all-width occupation bound, but bounded width at every positive error (Theorem~\ref{thm:adiclinear55}).

The quantitative theory has separate arithmetic hypotheses. For $s$ planar rotations with a dual badly approximable angle vector, the width is $\Theta(N^{s/(2s+1)})$ at every fixed error below half the signal amplitude when $0<\rho<1$. At $\rho=1$ the same law holds at each positive subcritical error, by Corollary~\ref{cor:unitnoise55}; the exact endpoint is not inferred by a limiting argument. Neither rational matrix entries nor an arbitrary enlarged alphabet implies the hypotheses of the matching upper bound.

The matrix-group structure used in purification is classical; in particular, compact groups of nonnegative matrices and their permutation structure were studied by Brown, Schwarz and Flor \cite{Flor}. Our reduction applies that structure to the joint closure of physical products and stochastic rows and preserves every seed--query error constraint. We give the necessary argument rather than assume that group relations already hold on hidden states. The remaining proofs use Haar projection, representative functions, constrained radius localization, executable word packing, and polygon enclosure. Section~\ref{sec:comparison54} separates these ingredients from the realization statements established here.
